% =====================================================================
%  Chapter 8 — Principle: Simplicity (One Way to Do Each Thing)
%  Map: research/chapter-map.md (Chapter 8)
%  Shape: idea → origin → cost → what it opens up. Ideas, not how-to.
%  Rule: the author is never quoted in body text; short quotes only in
%  footnotes. Full statements live in interviews/.
% =====================================================================
\chaptersubtitle{One Way to Do Each Thing}
\chapter{Simplicity}
\label{ch:oneway}

Most programming languages offer a menu. Several ways to repeat something,
several ways to make a choice, several ways to name and rename things, and
machinery to fill in whatever the programmer left unsaid. Each option is
reasonable on its own. Together they mean that any task can be written many
ways, and every reader must recognise all of them. Silica offers one way to do
each thing, and puts the variety in how those few pieces are combined.

\section{One way, many combinations}

Silica has a single way to repeat, a single way to choose between
alternatives, and a single way to give a value a name. When a program chooses
between cases, it must account for every case, and the compiler checks that it
has. The point is not austerity for its own sake. It is that a reader---human or
machine---should never have to wonder which of several near-identical forms the
author meant, or why.

\section{Show the shape, not a name}

The boldest refusal concerns how data is described. Most languages let
programmers invent names for kinds of data, and then names for those names.
Silica does not. The shape of the data is written out wherever it is used, so a
reader always sees what a thing actually is rather than a label that must be
looked up somewhere else. Silica also refuses to convert silently between kinds
of data; a value is produced as the kind it needs to be.

This was not always so. The ability to name types existed in Silica's first
months and was removed. The reason came from watching language models work.
During the months when models were writing much of Barney's code
(Chapter~\ref{ch:why}), the thing they struggled with most was named types of
any real depth: names defined in terms of other names, defined in terms of still
others. Barney saw that deep inheritance trees in object-oriented languages
would trap them the same way. A shape written out in full gives a model, and a
person, nothing to unwind.\footnote{Interview, 2026-11-13.}

Generic types, a feature most modern languages treat as essential, went six days
after the first commit, with a note that said it was for now. It is now a
conviction. Generics were meant to make computation less complicated; in
Barney's experience, once theory met the real world, they did not. They became
something closer to a fashion, so that a language without them was assumed not
to be a real language. What helps a programmer reason realistically about code
is being able to see the shape of the data being worked on.

Silica does have traits, which let different kinds of data share behaviour.
Barney regards them as a different idea from generics, and an older one: traits
existed before both generics and object-oriented classes, and they feel
different enough that newcomers do not mistake one for the
other.\footnote{Interview, 2026-11-13.}

% TODO-INTERVIEW: The standard library's collections are told what kinds of
%   values they hold. Clarify for readers how that differs from generics.

\section{Written for people and for machines}

From its founding notes, Silica was meant to be easy for language models, as
well as people, to read and write. That aim did not grow out of building Silica
with AI agents; it came first. It began with the experience of relying on models
to write code during a recovery, and seeing them fight languages that had been
designed to be understood by humans alone (Chapter~\ref{ch:why}). Every refusal
in this chapter removes a place where a generator could produce something that
looks right and is not. When there is one way to say a thing, there is one thing
to check.

The project's beginner book draws the conclusion boldly. It advises newcomers
not to type programs by hand at all, but to let a language model write them,
and to spend their own effort learning to read and judge what comes
back.\footnote{\src{silica/docs/learn-programming.md}, line 15.} Yet the
project holds people to the opposite standard in human conversation: every
contribution belongs to the person who submits it, and no one may present a
model's words as their own. The machine types; the human judges, and answers
for the result.

\section{The temperament behind it}

The same instinct shows outside the language. A small tool the project designed
for its own work was defined first by everything it would not be: not a
full editor, not a platform, not dependent on someone else's kit. Small,
native, understandable. Silica's simplicity is that temperament applied to a
language.

% TODO-INTERVIEW: Where does that instinct come from for you?

\section{What it costs}

Some Silica programs are longer than their equivalents elsewhere, because they
spell out what other languages abbreviate. Experienced programmers miss
favourite features. And some patterns that are one line in another language take
a moment's thought to express with fewer, plainer pieces.

\section{What it opens up}

\begin{itemize}
  \item \textbf{Machine-written code that people can trust.} When a language
        model has fewer ways to be almost right, its output is more often
        simply right, and easier to check when it is not.
  \item \textbf{A new way to learn programming.} Beginners can start by
        reading and judging programs rather than by memorising syntax, because
        there is less syntax to memorise and less ambiguity to decode.
  \item \textbf{Code anyone can pick up.} With one way to do each thing, a
        program written by a stranger reads like one written by a colleague.
  \item \textbf{Simpler tools.} Analysers, editors and compilers have less
        to understand, so they can understand it more deeply.
\end{itemize}

% TODO-INTERVIEW: What does a world of human judges and machine typists look
%   like to you, ten years out?
